\icode{E-MI-SI-L3-261113-001}
如图所示，\tikz\draw (0,0)--(60:.36)--(.18,0);\;形硬钢丝竖直固定放置，$OM$ 沿竖直方向，$\angle MON =$ \,\qty{30}{\degree}，一个光滑的轻环套在 $ON$ 上，一根足够长的轻绳穿过轻环，一端固定在 $A$ 点，另一端悬挂质量为 $m$ 的物体。当系统处于稳定状态时，$ON$ 受到轻环的弹力大小为\paren[]

\begin{tikzpicture}[font=\scriptsize]
	\draw[thick] (0,0)node[left]{\(N\)}--++(60:3)coordinate(O)node[right]{\(O\)}-|(O|-0,0)node[right]{\(M\)};
	\draw[magenta] (60:1.5)coordinate(a)--++(30:{sqrt(3)/2})node[right]{\(A\)} (a)--++(0,-.8)coordinate(b);
	\draw[fill=orange] ($(b)+(-.1,0)$) rectangle ($(b)+(.1,-.2)$);
	\draw[rotate=60] (1.5,0) ellipse(.03 and .06);
	\draw ($(O)+(0,-.3)$) arc(-90:-120:.3);
\end{tikzpicture}
\begin{choices}
	\item $\frac{1}{2}mg$
	\item $\frac{\sqrt{3}}{2}mg$
	\item $mg$
	\item $\sqrt{3}mg$
\end{choices}
\begin{answer}
	C\qquad 提示：
\end{answer}